% ! TEX root = ../m3-20-21.tex

\chapter{Transformarea Fourier}

Transformarea Fourier, cu diversele sale variante (discretă, rapidă etc.)
este foarte utilă pentru studiul undelor sinusoidale. Dacă o \emph{serie}
Fourier dezvoltă o funcție într-o serie în care termenii sînt sinusuri
și cosinusuri, \emph{transformata} Fourier este o transformare integrală,
în care aspectele periodice corespunzătoare funcțiilor trigonometrice
se păstrează.

De aceea, veți întîlni foarte des transformatele Fourier în cazul analizelor
semnalelor sonore, fie pentru digitalizare sau pentru diverse descompuneri.
Un exemplu concret este în inteligența artificială și învățarea automată,
unde se lucrează la programe care să facă recunoaștere vocală, transcrieri
și traduceri. În toate aceste aplicații, o analiză a semnalului audio cu
ajutorul transformatei Fourier este indispensabil.

%%%%%%%%%%%%%%%%%%%%%%%%%%%%%%%%%%%%%%%%%%%%%%%%%%%%%%%%%%%%%%%%%%%%%%
\section{Elemente teoretice}

Foarte pe scurt, vom da definițiile și principalele proprietăți care vor
fi de folos în exerciții.

Avem nevoie de următoarea noțiune:
\begin{definition}\label{def:functie-l1}\index{funcție!$L^1$}
    Fie $ f : \RR \to \RR $ o funcție reală. Spunem că ea este $ L^1 $
    (formal, \emph{aparține mulțimii $ L^1(\RR) $}), dacă are proprietatea:
    \[
        \int_{-\infty}^\infty |f(t)| dt < \infty.
    \]
    Intuitiv, mutînd tot graficul funcției deasupra axei $ Ox $ (prin
    oglindire pe intervalele negative), aria de sun grafic este finită.
    Deci funcția nu are nicio zonă în care \enquote{explodează} spre
    $ \pm \infty $ asimptotic.
\end{definition}

Acestea sînt funcțiile cărora le vom defini transformările Fourier. Deci,
dacă nu se precizează altfel, toate funcțiile care se transformă vor fi
presupuse $ L^1 $.

Definiția principală urmează.

\begin{definition}\label{def:transf-f}\index{transformata!Fourier}
    Fie $ f $ o funcție $ L^1 $. Se definește \emph{transformata Fourier} a lui
    $ f $, notată $ \kal{F} $ (alternativ, $ \kal{F}[f] $ sau, și mai
    explicit, $ \kal{F}[f(t)](\omega) $)\footnotemark funcția complexă
    definită prin:
    \[
        \kal{F}[f] : \RR \to \CC, \quad \kal{F}[f](\omega) = %
        \int_{-\infty}^\infty f(t) e^{-i\omega t} dt.
    \]
\end{definition}

\footnotetext{Pentru a nu încărca inutil notația, vom folosi această
    variantă explicită doar aici. Regula care se va păstra în continuare
    este următoarea: argumentul funcției $ f $ care se transformă este $ t $,
    iar argumentul transformatei $\kal{F}[f] $ este $ \omega $.
    De asemenea, argumentul $ t $ al funcției $ f $ se mai numește \emph{timp}, iar
argumentul $ \omega $ al transformatei se mai numește \emph{frecvență}.}

Se vede aici importanța faptului că $ f $ este $ L^1 $.

Folosind formula lui Euler pentru exponențiala complexă care apare în
integrală, precum și paritatea funcțiilor trigonometrice, facem următoarele
observații imediate:
\begin{itemize}
    \item Dacă $ f $ este funcție pară, atunci transformata Fourier are partea
        imaginară nulă (integrala sinsului pe un interval simetric față de origine ---
        în acest caz special, $ (-\infty, \infty) $ --- este nulă, iar integrala
        funcției pare cosinus este dublul integralei calculate pe jumătate din
        interval) și rămînem cu:
        \[
            \kal{F}[f](\omega) = 2 \int_0^\infty f(t)\cos(\omega t) dt;
        \]
    \item Dacă $ f $ este funcție impară, folosind argumentul din cazul anterior,
        transformarea Fourier are partea reală nulă și rămînem cu:
        \[
            \kal{F}[f](\omega) = -2i \int_0^\infty f(t) \sin (\omega t) dt.
        \]
\end{itemize}

Cele două cazuri speciale se numesc, respectiv, \emph{transformarea (Fourier) prin sinus}
și \emph{transformarea (Fourier) prin cosinus} ale funcției $ f $.

Transformarea Fourier se poate inversa relativ simplu:
\begin{theorem}[Transformarea Fourier inversă]\label{thm:fourier-inv}
    \index{transformata!Fourier!inversă}
    Fie $ f $ o funcție $ L^1 $ și $ \kal{F}[f] $ transformata sa Fourier.
    Presupunem că și $ \kal{F}[f] $ este funcție $ L^1 $ și atunci
    se poate recupera funcția $ f $ prin formula:
    \[
        f(t) = \dfrac{1}{2 \pi} \int_{-\infty}^\infty \kal{F}[f]e^{i\omega t}d\omega.
    \]
\end{theorem}

Dacă funcția $ f $ care se transformă are \enquote{puncte cu probleme}
(care sigur aparțin domeniului de integrare, deoarece acesta este întreg
$ \RR $), atunci transformarea Fourier se calculează cu ajutorul reziduurilor.
Fie, așadar, $ p_1, \dots, p_n $ poli ai funcției $ f $ și presupunem că
$ \ds\lim_{|z| \to \infty} f(z) = 0 $. Atunci:
\begin{itemize}
    \item Dacă $ \dr{Im}p_k > 0 $, atunci:
        \[
            \kal{F}[f](\omega) = 2\pi i \sum_{k=1}^n \dr{Rez}\left(f(z)e^{-i\omega z}, p_k\right), %
            \quad \text{ pentru } \omega < 0;
        \]
    \item Dacă $ \dr{Im}p_k < 0 $, atunci:
        \[
            \kal{F}[f](\omega) = -2\pi i \sum_{k=1}^n \dr{Rez}\left(f(z)e^{-i\omega z}, p_k\right), %
            \quad \text{ pentru } \omega > 0;
        \]
\end{itemize}

În exerciții, ne vor mai fi de folos și următoarele proprietăți ale
transformatei Fourier:
\begin{enumerate}[(1)]
    \item \textbf{Liniaritatea:} $ \kal{F}[\alpha f + \beta g] = \alpha \kal{F}[f] + \beta \kal{F}[g] $,
        pentru orice $ f, g $ funcții $ L^1 $ și $ \alpha, \beta \in \CC $ scalari;
    \item \textbf{Simetria:} $ \kal{F}\left[\kal{F}[f](\omega)\right] = 2\pi f(-\omega) $;
    \item \textbf{Rescalarea:} Fie $ \alpha \in \CC $. Atunci:
        \[
            \kal{F}\left[f\left(\alpha t\right)\right](\omega) = %
            \dfrac{1}{|\alpha|}\kal{F}[f]\left(\dfrac{\omega}{\alpha}\right).
        \]
    \item \textbf{Translația după $ t $:}
        \[
            \kal{F}[f(t - t_0)] = \kal{F}[f(t)] \cdot e^{-i\omega t_0}, %
            \quad \forall t_0 \in \RR;
        \]
    \item \textbf{Translația după $ \omega $:}
        \[
            \kal{F}\left[e^{i\omega_0 t}f(t)\right] = \kal{F}[f(t)](\omega - \omega_0), %
            \quad \forall \omega_0 \in \RR;
        \]
    \item \textbf{Derivarea după $ t $:}
        \[
            \kal{F}[f^{(n)}(t)] = (i\omega)^n \kal{F}[f](\omega),
        \]
        în ipoteza că $ f $ este de $ n $ ori derivabilă, pentru un anume $ n $;
    \item \textbf{Derivarea după $ \omega $:}
        \[
            \kal{F}\left[(-it)^n f(t)\right] = \kal{F}^{(n)}[f](\omega),
        \]
        în ipoteza că derivatele de pînă la $ n $ ori au sens;
    \item \textbf{Transformata conjugatei complexe:} Notăm $ f^\ast(t) $ funcția
        conjugată complexă asociată lui $ f(t) $. Atunci:
        \[
            \kal{F}[f^\ast(t)] = \left(\kal{F}[f(t)]\right)^\ast(-\omega).
        \]
        Altfel spus, se transformă funcția inițială, apoi se ia conjugatul
        rezultatului și se schimbă semnul argumentului;
    \item \textbf{Convoluția în timp:} Definim:
        \[
            h(t) = (f \ast g)(t) = \int_{-\infty}^\infty f(\tau)g(t - \tau)d\tau
        \]
        produsul de convoluție al funcțiilor $ f $ și $ g $.
        Atunci:
        \[
            \kal{F}[f \ast g] = \kal{F}[f] \cdot \kal{F}[g];
        \]
    \item \textbf{Convoluția în frecvență:}
        \[
            \kal{F}\left[f(t) \cdot g(t)\right] = %
            \dfrac{1}{2 \pi} \int_{-\infty}^\infty \kal{F}[f](y) \cdot \kal{F}[g](\omega - y)dy.
        \]
\end{enumerate}

%%%%%%%%%%%%%%%%%%%%%%%%%%%%%%%%%%%%%%%%%%%%%%%%%%%%%%%%%%%%%%%%%%%%%%
\section{Aplicații la ecuații integrale}

Folosind formula de transformare Fourier, precum și inversarea acesteia,
putem rezolva ecuații integrale. Vom prezenta acest lucru pe un exemplu.

\textbf{Exemplu:} Rezolvăm ecuația integrală:
\[
    \int_0^\infty y(t) \cos tx dt = \dfrac{1}{x^2 + 1}.
\]

\emph{Soluție:} Pentru a face legătura cu transformările Fourier,
avem nevoie să prelungim funcția $ y $ (astfel încît integrala
să fie pe $ (-\infty, \infty) $). Deoarece integrandul este par,
acest lucru conduce la dublarea rezultatului. Deci avem:
\[
    \int_{-\infty}^\infty y(t) \cos tx dt = \dfrac{2}{x^2 + 1}, \quad %
    \int_{-\infty}^\infty y(t) \sin tx dt = 0.
\]

Acest lucru conduce la:
\[
    \dfrac{1}{ 2\pi} \int_{-\infty}^\infty y(t) e^{-itx} dt = \dfrac{1}{\pi(x^2 + 1)}.
\]

Recunoaștem în membrul stîng formula de trasnformare Fourier,
deci rezultă:
\[
    \kal{F}[y(t)](x) = \dfrac{1}{\pi(x^2 + 1)}.
\]

Folosim acum formula de inversare și rezultă:
\[
    y(t) = \dfrac{1}{\pi} \int_{-\infty}^\infty \dfrac{1}{x^2 + 1} e^{itx} dx.
\]

Acest calcul se poate face cu teorema reziduurilor, deoarece integrandul
are doi poli, $ p_{1,2} = \pm i $, ambii simpli. Obținem:
\begin{itemize}
    \item Pentru $ t < 0 $ și polul $ p_1 = i $:
        \begin{align*}
            y(t) &= \dfrac{1}{\pi} \cdot 2 \pi i \dr{Rez}\left(\dfrac{1}{x^2 + 1} e^{itx}, i\right) \\
                 &= 2 i \cdot \lim_{z \to i} (z - i) \dfrac{1}{z^2 + 1} e^{itz} \\
                 &= 2i \cdot e^{-t};
        \end{align*}
    \item Pentru $ t > 0 $ și polul $ p_2 = -i $:
        \begin{align*}
            y(t) &= \dfrac{1}{pi} \cdot - 2\pi i \dr{Rez}\left( \dfrac{1}{x^2 + 1} e^{itx}, -i\right) \\
                 &= -2i \cdot \lim_{z \to -i} (z + i) \dfrac{1}{z^2 + 1} e^{itz} \\
                 &= 2i e^{t}.
        \end{align*}
\end{itemize}

Concluzia este că soluția arată astfel:
\[
    y(t) = \begin{cases}
        2i \cdot \exp(-t), & t < 0 \\
        2i \cdot \exp(t), & t > 0
    \end{cases} = 2i \cdot \exp\left( |t| \right).  
\]

%%%%%%%%%%%%%%%%%%%%%%%%%%%%%%%%%%%%%%%%%%%%%%%%%%%%%%%%%%%%%%%%%%%%%%

\section{Exerciții}

1. Rezolvați ecuațiile integrale:
\begin{enumerate}[(a)]
    \item $ \ds\int_0^\infty f(t) \sin tx dt = e^{-x}, \quad x > 0 $;
    \item $ \ds\int_0^\infty f(t) \cos tx dt = \dfrac{1}{(1 + x^2)^2} $;
\end{enumerate}

\vspace{0.3cm}

2. Calculați transformatele Fourier pentru funcțiile:
\begin{enumerate}[(a)]
    \item $ f(t) = \begin{cases} 1 - |t|, & t \in [-1, 1] \\ 0, & \text{ în rest} \end{cases} $;
    \item $ f(t) = \dfrac{t}{t^2 + a^2}, a > 0 $;
    \item $ f(t) = \dfrac{1}{(t^2 + 1)^2} $;
    \item $ f(t) = \begin{cases} \exp(2t), & t \leq 0 \\ 0, & t > 0 \end{cases} $;
    \item $ f(t) = \begin{cases} t^2 - t, & 0 < t < 1 \\ 0, & \text{ în rest} \end{cases} $;
    \item $ f(t) = \begin{cases} t, & |t| < 1 \\ 0, & |t| > 1 \end{cases} $.
\end{enumerate}

\vspace{0.3cm}

3. Calculați transformatele Fourier pentru funcțiile de mai jos,
folosind teorema reziduurilor:
\begin{enumerate}[(a)]
    \item $ f(t) = \dfrac{2}{t^2 + 9} $;
    \item $ f(t) = \dfrac{t}{(t^2 + 9)(t^2 + 1)} $;
    \item $ f(t) = \dfrac{1}{(t^2 + 4)^2} $.
\end{enumerate}
